% ============================================================
%  H. Høffding, Psykologi i Omrids paa Grundlag af Erfaring
%  2. til dels ændrede Udgave.
%  København: P. G. Philipsens Forlag, 1885. IV + 421 s.
%  TRANSCRIPTION (Danish)
%  Source: KB scan (pdf).
% ============================================================
\documentclass[12pt,a4paper]{book}
\usepackage[utf8]{inputenc}
\usepackage[T1]{fontenc}
\usepackage[danish]{babel}
\usepackage{microtype}
\usepackage{setspace}
\usepackage[margin=1.2in]{geometry}
\usepackage{fancyhdr}
\setlength{\headheight}{28pt}
\addtolength{\topmargin}{-13.5pt}
\usepackage{hyperref}

\hypersetup{
  pdftitle={Psykologi i Omrids},
  pdfauthor={Harald Høffding},
  hidelinks
}

\pagestyle{fancy}
\fancyhf{}
\fancyhead[LE,RO]{\thepage}
\fancyhead[RE]{\textit{Psykologi i Omrids}}
\fancyhead[LO]{\textit{\leftmark}}

\setstretch{1.3}

\title{\textbf{Psykologi i Omrids}\\[6pt]
  \large paa Grundlag af Erfaring\\[6pt]
  \normalsize Anden, til dels ændrede Udgave}
\author{Harald Høffding}
\date{København: P.\ G.\ Philipsens Forlag, 1885}

\begin{document}
\maketitle
\thispagestyle{empty}
\newpage

\tableofcontents
\newpage


%% ============================================================
%%  I. PSYKOLOGIENS GENSTAND OG METODE (pp. 1--37)
%% ============================================================
\chapter{Psykologiens Genstand og Metode}
\label{ch:genstand}

% 1. Foreløbig Betegnelse af Psykologien som Lære om Sjælen.
% 2. Ydre Iagttagelse gaar forud for indre.
% 3. Sprogets Vidnesbyrd.
% 4. Psykologisk Udvikling af Forskellen mellem Jeg og Ikke-Jeg.
% 5. Den mytologiske Opfattelse af Sjælen.
% 6. Direkte og indirekte Opfattelse af Sjæleliv.
% 7. Psykologi og Metafysik.
% 8. Psykologiens Metode.
%    a. Indre Iagttagelsers Vanskelighed.
%    b. Individuelle Forskelligheders Indflydelse.
%    c. Psykologisk Analyse.
%    d. Experimental Psykologi.
%    e. Subjektiv og objektiv Psykologi.
%    f. De forskellige Metoders indbyrdes Forhold.
% 9. Psykologiens Forhold til Logik og Etik.

% [text to be added]


%% ============================================================
%%  II. SJÆL OG LEGEME (pp. 38--80)
%% ============================================================
\chapter{Sjæl og Legeme}
\label{ch:sjael-legeme}

% 1. Empirisk (fænomenologisk) Standpunkt.
% 2. Loven om Energiens Bestaaen.
% 3. Det organiske Liv og Energiens Bestaaen.
% 4. a. Nervesystemets Betydning.
%    b. Reflexbevægelse.
%    c. Underordnede Nervecentra.
%    d. Den store Hjerne.
%    e. Den store Hjerne og de lavere Centra.
% 5. Foreløbig Karakteristik af Bevidsthedslivet.
% 6. Paralleler mellem Bevidstheden og Nervesystemet.
% 7. Proportionalitet mellem Bevidsthedsliv og Hjernevirksomhed.
% 8. a. Dualistisk-spiritualistisk Hypotese.
%    b. Monistisk-materialistisk Hypotese.
%    c. Monistisk-spiritualistisk Hypotese.
%    d. Identitetshypotesen (Monismen).

% [text to be added]


%% ============================================================
%%  III. DET BEVIDSTE OG DET UBEVIDSTE (pp. 81--98)
%% ============================================================
\chapter{Det bevidste og det ubevidste}
\label{ch:bevidst-ubevidst}

% 1. Definition af det ubevidste.
% 2. Bevidst Tankeresultat af ubevidst Forarbejde.
% 3. Bevidst Sansningsresultat af ubevidst Forarbejde.
% 4. Ubevidste Mellemled.
% 5. Instinkt og Vane.
% 6. Ubevidst Virksomhed samtidig med bevidst.
% 7. Følelsens ubevidste Væxt.
% 8. Drømmetilstanden.
% 9. Opvaagnen ved Indtrykkets psykiske Relation.
% 10. Hypotese om Sjælelivets Udstrækning.
% 11. Psykologi og fysisk Mekanik.
% 12. Fælles Love for Aand og Materie.

% [text to be added]


%% ============================================================
%%  IV. INDDELING AF DE PSYKOLOGISKE ELEMENTER (pp. 99--114)
%% ============================================================
\chapter{Inddeling af de psykologiske Elementer}
\label{ch:inddeling}

% 1. Inddeling af Elementer, ej af Tilstande.
% 2. Den psykologiske Tredeling.
% 3. Tredelingen ikke oprindelig.
% 4. Den individuelle Bevidstheds Udvikling.
% 5. Psykologisk Differentiation under Slægtens Udvikling.
% 6. Differentiationens Betingelser.
% 7. a. Ingen Erkendelse uden Følelse.
%    b. Ingen Erkendelse uden Villen.
%    c. Ingen Følelse uden Erkendelse.
%    d. Sammenhæng mellem Følelse og Villie.
%    e. Villien som det første og det sidste.

% [text to be added]


%% ============================================================
%%  V. ERKENDELSENS PSYKOLOGI (pp. 115--254)
%% ============================================================
\chapter{Erkendelsens Psykologi}
\label{ch:erkendelse}

\section{Fornemmelse}
\label{sec:fornemmelse}

% 1. Den psykologiske Betydning af Spørgsmaalet om
%    Fornemmelsernes Enkelthed og Selvstændighed.
% 2. Fornemmelsernes Enkelthed.
% 3. Fornemmelsernes Selvstændighed.
% 4. Om Fornemmelsernes Kvalitet.
% 5. Forholdsloven paa Fornemmelsernes Omraade.
% 6. Bevægelsesfornemmelser.
% 7. Sansning og Bevægelse.

% [text to be added]

\section{Forestilling}
\label{sec:forestilling}

% 1. Fornemmelse og Iagttagelse.
% 2. Frie Forestillinger.
% 3. Fornemmelse, Perception og fri Forestillen.
% 4. Frie Forestillingers Udsondring fra Iagttagelser.
% 5. Formal og real Enhed i Bevidstheden.
% 6. Forestillingernes Bevarelse.
% 7. a. Erindringsbilleder, Hallucinationer og Illusioner.
%    b. Erindring betinget ved Oplevelsesvilkaarene.
%    c. Erindring betinget ved Genkaldelsesvilkaarene.
%    d. Erindring betinget ved Forestillingernes Beskaffenhed.
% 8. a. Om Forestillingsforbindelsens Lovbundethed.
%    b. Forestillingsforbindelsens Love.
%    c. Grundlov for Forestillingsforbindelserne.
%    d. Forglemmelsens Love.
% 9. Enkeltforestillinger, Individualforestillinger og
%    Almenforestillinger.
% 10. Sproget og Forestillingerne.
% 11. Idéassociation og Tænkning.
% 12. Dannelsen af frie konkrete Individualforestillinger (Fantasi).

% [text to be added]

\section{Tids- og Rumsopfattelse}
\label{sec:tid-rum}

% 1--10. [subsections as in TOC]

% [text to be added]

\section{Opfattelsen af det virkelige}
\label{sec:virkelighed}

% 1. Erkendelsesindholdet som Udtryk for en Virkelighed.
% 2. Sammenhæng som Virkelighedskriterium.
% 3. Aarsagsforholdet.
% 4. Aarsagsbegrebets psykologiske Udvikling.
% 5. Erkendelsens Grænser.

% [text to be added]


%% ============================================================
%%  VI. FØLELSENS PSYKOLOGI (pp. 255--356)
%% ============================================================
\chapter{Følelsens Psykologi}
\label{ch:foelelse}

\section{Følelse og Sansefornemmelse}
\label{sec:foelelse-sanse}

% 1. Følelseslivets Enhed.
% 2. Følelse forskellig fra Sansefornemmelse.
% 3. Følelsen og de enkelte Sanser.
%    a. Livsfølelsen.  b. Følelser ved Berøring og Bevægelse.
%    c. Følelser ved Smag.  d. Følelser ved Lugt.
%    e. Følelser ved Syn og Hørelse.
% 4. Elementarfølelsernes naturlige Udviklingsgang.

% [text to be added]

\section{Følelse og Forestilling}
\label{sec:foelelse-forestilling}

% 1. Følelsens Oprindelighed.
% 2. a. Afsky, Sorg, Had.  b. Kærlighed, Glæde, Sympati.
%    c. Drift, Begær.  d. Haab, Frygt.  e. Blandede Følelser.
% 3. Loven for Følelsens Udvikling.
% 4. Erindren af Følelser.

% [text to be added]

\section{Egoistisk og sympatisk Følelse}
\label{sec:egoistisk-sympatisk}

% 1--9 [subsections as in TOC]

% [text to be added]

\section{Følelsens Fysiologi og Biologi}
\label{sec:foelelse-fysiologi}

% 1--3 [subsections as in TOC]

% [text to be added]

\section{Forholdslovens Gyldighed for Følelserne}
\label{sec:forholdslov-foelelse}

% 1--9 [subsections as in TOC, including det ophøjede and det latterlige]

% [text to be added]

\section{Følelsens Indflydelse paa Erkendelsen}
\label{sec:foelelse-erkendelse}

% 1--4 [subsections as in TOC]

% [text to be added]


%% ============================================================
%%  VII. VILLIENS PSYKOLOGI (pp. 357--421)
%% ============================================================
\chapter{Villiens Psykologi}
\label{ch:villie}

\section{Villiens Oprindelighed}
\label{sec:villie-oprind}

% 1. Villien den mest primitive og den mest afledede sjælelige Ytring.
% 2. Spontaneitet og Irritabilitet.
% 3. Spontane og reflexe Bevægelser hos højere Organismer.
% 4. Instinktets og Villiens fysiologiske Sæde.
% 5. Uvilkaarlig og vilkaarlig Opmærksomhed.
% 6. a. Villien og Bevægelsesforestillingerne.
%    b. Isolation og Kombination af Bevægelser.
%    c. Det medfødte Grundlags Betydning.

% [text to be added]

\section{Villien og de andre Bevidsthedselementer}
\label{sec:villie-bevidsthed}

\subsection{Villiens højere Udvikling, betinget ved Erkendelsens og
  Følelsens Udvikling}
\label{ssec:villie-hoiere}

% a. Driftens Psykologi.
% b. Ønsket.
% c. Overvejelse, Forsæt og Beslutning.

% [text to be added]

\subsection{Villiens Tilbagevirken paa Erkendelse og Følelse}
\label{ssec:villie-tilbagevirken}

% a. Villiens Tilbagevirken paa Erkendelsen.
% b. Villiens Tilbagevirken paa Følelsen.
% c. Villiens Tilbagevirken paa sig selv.

% [text to be added]

\subsection{Modsætningsforhold mellem Villien og de andre
  Bevidsthedselementer (Koncentration og Expansion)}
\label{ssec:koncentration-expansion}

% [text to be added]

\subsection{Bevidstheden om Villien}
\label{ssec:bevidsthed-villie}

% a. Bevidstheden om at ville.
% b. Virkelighedsproblemet paa Villiens Omraade.

% [text to be added]

\subsection{Villien og det ubevidste Sjæleliv}
\label{ssec:villie-ubevidst}

% a. Bevidsthedens Midtpunkt ikke altid Individualitetens Midtpunkt.
% b. Determinisme og Indeterminisme.

% [text to be added]

\section{Den individuelle Karakter}
\label{sec:karakter}

% 1. Typiske individuelle Forskelligheder.
% 2. Fysiske, sociale og nedarvede Elementer.
% 3. Psykisk Individualitet en faktisk Grænse for Erkendelsen.
% 4. Psykologien og Udviklingshypotesen.

% [text to be added]

\end{document}
